% !PRINTED 37
\weberpage{37}
From (12) it follows further that
\[
1 = \Theta^a = A^{\frac{\lambda a}{b}};
\]
hence $\lambda a : b$ must be a multiple of $a$, and consequently $\lambda$ a multiple of $b$. If we therefore set
\weq{13}{B = \Theta A^{-\frac{\lambda}{b}}}
then $B^b = 1$, and at the same time $B^b$ is the lowest power of $B$ that becomes equal to a power of $A$, because, if $B^{b'}$ is a power of $A$, the same, by (13), holds also of $\Theta^{b'}$. We assume the element $\Theta$ to be so chosen that $b$ becomes as great as possible. Then for every other element $\Theta_1$ of $S$ as well, $\Theta_1^b$ is always a power of $A$ whose exponent is divisible by $b$.

If now
\[
\begin{gathered}
\alpha = 0, 1, 2 \ldots a-1 \\
\beta = 0, 1, 2 \ldots b-1,
\end{gathered}
\]
then the elements
\weq{14}{A^\alpha B^\beta}
are, $ab$ in number, all different from one another. They form a divisor $S'$ contained in $S$, which is represented by the basis $A, B$. If $S'$ is identical with $S$, we have reached our goal.

But if $S$ is not yet exhausted by (14), we continue in the same manner.

By an application of complete induction we shall draw the conclusion at once in full generality.

Let a divisor $S_{\nu-1}$ of $S$ have been found, represented by a basis in the form
\weq{15}{S_{\nu-1} = A_1^{\alpha_1} A_2^{\alpha_2} \ldots A_{\nu-1}^{\alpha_{\nu-1}},}
about which we make the following assumptions:

1) Let the degrees of $A_1, A_2 \ldots A_{\nu-1}$ be $a_1, a_2 \ldots a_{\nu-1}$, and let
\[
a_1 \geqq a_2 \geqq \ldots \geqq a_{\nu-1}.
\]

2) For every element $\Theta$ contained in $S$ let
\weq{16}{\Theta^{a_{\nu-1}} = A_1^{\lambda_1} A_2^{\lambda_2} \ldots A_{\nu-2}^{\lambda_{\nu-2}},}
that is, contained in $S_{\nu-2}$, and let the exponents $\lambda_1, \lambda_2 \ldots \lambda_{\nu-2}$ be divisible by $a_{\nu-1}$.

The group $S'$ derived above satisfies these requirements if $\nu = 3$ and $a_{\nu-1} = b$ is set, and it is now to be shown
%CONT>>
